\documentclass[10pt]{article}
\usepackage[margin=20mm]{geometry}
\usepackage{amsmath,amssymb,mathtools,fontspec,unicode-math,xurl,hyperref}
\IfFontExistsTF{DejaVu Serif}{\setmainfont{DejaVu Serif}}{\setmainfont{Latin Modern Roman}}
\setmathfont{latinmodern-math.otf}
\setlength{\parindent}{0pt}\setlength{\parskip}{6pt}\sloppy
\title{Verification of the quantitative Struwe decomposition}
\author{AI-assisted technical review}
\date{}
\begin{document}
\maketitle
Abstract: A claim-focused examination of the author manuscript identifies three substantive issues. The unrestricted six-dimensional distance estimate in Corollary 1.6 fails for an explicit positive multiple of one bubble whose residual is exactly one. The decomposition used in Lemma 6.2 does not have the asserted orthogonality and therefore does not justify its spectral-gap argument. The final unweighted estimate in Lemma 6.5 does not give a little-oh interaction remainder in dimensions at least twelve. A conditional globalization correction and a local weighted replacement for the last estimate are provided. The findings concern the identified author manuscript and the stated proof steps; they do not establish that the local sharp stability theorem is false.\par
\section*{Identity, reviewed material, and scope}
This report concerns Sharp quantitative estimates of Struwe’s decomposition, by Bin Deng, Liming Sun, and Jun-cheng Wei, associated with DOI 10.1215/00127094-2024-0026 (https://doi.org/10.1215/00127094-2024-0026). The primary source is the 59-page author-hosted manuscript dated January 3, 2024 [DSW24]. All source page numbers below are one-based PDF pages, which agree with the printed pages. The publisher’s version of record was not acquired or compared. The preserved arXiv v2 source is dated 2021 and is a separate version [DSW21].\par
The review checks the statements in the introduction, the closing argument of Section 2, the energy and remainder estimates in Section 6, the relevant sharpness argument in Section 7, and the weight bounds needed for a local correction. The lengthy bubble-tree and neck-region estimates in Sections 3–5 are not independently certified in full. No missing optional code is treated as a mathematical defect. The supported assessment is Major: the first finding disproves a stated quantitative corollary, and the other two identify invalid steps used in the proof of the central theorem.\par
\section*{Definitions and operative assumptions}
Equation reference: eq:rate. The manuscript uses \[p=\frac{n+2}{n-2},\qquad
U[z,\lambda](x)=(n(n-2))^{(n-2)/4}
 \left(\frac{\lambda}{1+\lambda^2|x-z|^2}\right)^{(n-2)/2}.\] These satisfy \(-\Delta U=U^p\) and have parameter-independent \(\dot H^1\) norm. The residual is \(\Gamma(u)=\|\Delta u+u|u|^{p-1}\|_{H^{-1}}\). We first use the dual of the homogeneous energy norm, consistent with the source’s \(\dot H^1\) framework. The rate in source equation (1.8), p. 4, is \[
 \zeta_n(t)=
 \begin{cases}t\sqrt{|\log t|},&n=6,\\
 t^{p/2},&n\ge7,\end{cases}\qquad t>0.\] Theorem 1.4 assumes proximity to a weakly interacting family of bubbles. Corollary 1.6 instead assumes positivity and an energy shell, without an explicit small-residual condition. Constants may depend on \(n\) and the fixed number \(\nu\) of bubbles. For \(\nu\ge2\), the interaction parameter is \[Q=\max_{i\ne j}\left(\frac{\lambda_i}{\lambda_j}+\frac{\lambda_j}{\lambda_i}
       +\lambda_i\lambda_j|z_i-z_j|^2\right)^{-(n-2)/2}.\] It tends to zero along the separated configurations considered below.\par
\section*{Substantive findings}
Problem 1 (The unrestricted six-dimensional corollary vanishes at a nonzero residual).  Corollary 1.6, p. 5, together with equation (1.8), p. 4, is false as written for \(n=6\).\par
\subsection*{Original assertion}
For every nonnegative \(u\in\dot H^1(\mathbb R^n)\) satisfying \[(\nu-\tfrac12)S^n\le\|\nabla u\|_2^2\le(\nu+\tfrac12)S^n,\] the corollary asserts a distance to a sum of \(\nu\) normalized bubbles at most \(C\zeta_n(\Gamma)\). Here \(S\) is normalized by \(S\|u\|_{2^*}\le\|\nabla u\|_2\). The proof on p. 13 consists of a reference to the proof of a corollary in the cited Figalli–Glaudo work.\par
\subsection*{Exact check}
Equation reference: eq:bubble. Set \(n=6\), \(\nu=1\), and \[
 U(x)=\frac{24}{(1+|x|^2)^2},\qquad
 A=\|\nabla U\|_2=\sqrt{\frac{1152\pi^3}{5}},\qquad
 c=\frac{1+\sqrt{1+4/A}}2,\qquad u=cU.\] The identity \(-\Delta U=U^2\) gives \[A^2=\int_{\mathbb R^6}U^3
 =24^3\pi^3\int_0^\infty\frac{r^5}{(1+r^2)^6}\,dr
 =\frac{1152\pi^3}{5}=S^6.\] Thus \(u>0\) and \(u\in\dot H^1\). Since \(A>5\) and \(c-1=1/(Ac)<1/5\), we have \(1/2<c^2<36/25<3/2\). Every energy hypothesis is satisfied.\par
Equation reference: eq:residual, eq:distance. Integration by parts and Cauchy–Schwarz show \(\|U^2\|_{\dot H^{-1}}\le A\), with equality on the test function \(U/A\). Consequently, \[
 \Gamma(u)=\|(c^2-c)U^2\|_{\dot H^{-1}}=Ac(c-1)=1.\] For every normalized bubble \(V=U[z,\lambda]\), \(\|V\|_{\dot H^1}=A\). The reverse triangle inequality and the choice \(V=U\) give the exact distance \[
 \inf_{z,\lambda}\|cU-U[z,\lambda]\|_{\dot H^1}=(c-1)A=\frac1c>0.\] The claimed right-hand side is \(C\zeta_6(1)=0\), a contradiction for every finite \(C\). The pair-interaction assertion is vacuous when \(\nu=1\).\par
\subsection*{Alternative readings and correction boundary}
Equation (1.8) expressly defines the rate for all positive arguments, and the source only says it is increasing near zero. The proximity assumption of Theorem 1.4 cannot silently be appended to the separately stated corollary. This example does not challenge the sufficiently small-residual asymptotic rate.\par
Even interpreting \(H^{-1}\) as the usual inhomogeneous dual does not remove the defect. Put \(B=\|U^2\|_{H^{-1}}\) and replace \(A\) by \(B\) in the definition of \(c\). Since \(\|U\|_2^2=96\pi^3\), testing against \(U\) gives \(B\ge A\sqrt{12/17}>5\). Then \(\Gamma=1\) and the same energy and positive-distance argument applies.\par
A precise local correction is to require \(0\le\Gamma\le\epsilon(n,\nu)<1\), sufficiently small for Struwe compactness to place the function in the local theorem’s regime. Alternatively replace the six-dimensional rate by \(\widetilde\zeta_6(t)=t\sqrt{1+|\log t|}\), with value zero at \(t=0\). This preserves the small-residual order. For \(t\ge\epsilon\) the modified rate is at least \(\epsilon\), while the energy shell bounds the distance to any fixed family of \(\nu\) bubbles. Increasing \(C\) then extends a valid local estimate to this remaining range; a fixed sufficiently separated family also bounds the pair interactions. This verifies the globalization mechanism conditional on the local theorem. The proof issues below must be addressed independently.\par
Problem 2 (The coefficients in the spectral decomposition are not projection coordinates).  Source equations (6.4)–(6.6), p. 46, do not imply the orthogonality used in the proof of Lemma 6.2 on p. 47.\par
\subsection*{Source step and complete algebraic check}
Equation reference: eq:gram. The source sets \[\rho_1=\sum_i\beta_iU_i+\rho_2,\qquad
 \beta_i=\langle \rho_1,U_i\rangle_{\dot H^1},\] and asserts \(\langle \rho_2,U_k\rangle_{\dot H^1}=0=\langle \rho_2,Z_k^a\rangle_{\dot H^1}\). Write \(G_{ki}=\langle U_i,U_k\rangle_{\dot H^1}\) and \(b_k=\langle \rho_1,U_k\rangle_{\dot H^1}\). The actual identity is \[
 (\langle \rho_2,U_k\rangle_{\dot H^1})_k=(I-G)b.\] The diagonal entries of \(G\) equal \(A_n^2=\|\nabla U\|_2^2\), not one, and its off-diagonal entries are positive interaction integrals. Furthermore, the initial tangent orthogonality of \(\rho_1\) only implies \[\langle \rho_2,Z_k^a\rangle_{\dot H^1}=-\sum_i b_i\langle U_i,Z_k^a\rangle_{\dot H^1},\] whose cross terms need not vanish.\par
Equation reference: eq:badgap. There is an exact arbitrarily small one-bubble diagnostic. Take \(n=6\), \(u=(1+t)U\), \(0<|t|\ll1\), with \(U\) as in [eq:bubble] (\#eq:bubble). The closest normalized bubble is \(U\), so \(\rho=tU\). The interaction forcing is zero and \(\rho_0=0\) solves the projected correction equation; hence \(\rho_1=tU\). Parameter invariance of \(\|U\|_{\dot H^1}\) gives \(\langle U,Z^a\rangle_{\dot H^1}=0\), as required. The printed definitions yield \[\beta=tA^2,\qquad
 \rho_2=t(1-A^2)U,\qquad
 \langle \rho_2,U\rangle_{\dot H^1}=t(1-A^2)A^2\ne0.\] This is an algebraic test of the asserted decomposition; it does not depend on assigning a value to a maximum over an empty set of bubble pairs. For this \(\rho_2\) the spectral quotient is exactly \[
 \frac{p\int U^{p-1}\rho_2^2}{\int|\nabla\rho_2|^2}=p=2.\] It therefore cannot satisfy the bound with \(\widetilde c<1\) displayed at the bottom of p. 47.\par
The problem also persists for arbitrarily weakly interacting multiple bubbles. In the orthogonal complement of all tangent modes, project \(U_1\) onto that complement and multiply it by a small scalar. As separation tends to infinity, the tangent Gram matrix tends to invertible single-bubble blocks and the cross terms tend to zero. For this admissible \(\rho_1\), the vector \(b\) tends to a nonzero multiple of the first coordinate vector, while \(G\to A_n^2I\). Equation [eq:gram] (\#eq:gram) is then nonzero whenever \(A_n^2\ne1\), including \(n=6\). Thus the diagnostic does not arise solely from the one-bubble endpoint.\par
\subsection*{Dependence, counterevidence, and repair}
Lemma 6.2 uses the asserted orthogonality both to eliminate multiplier terms and to invoke a second-variation gap. It feeds Lemma 6.3, Proposition 6.4 (p. 49), and the closing proof of Theorem 1.4 (p. 13). This is a core proof defect; equation [eq:badgap] (\#eq:badgap) does not disprove the conclusion of Lemma 6.2 or Theorem 1.4.\par
An intended genuine projection would resolve the elementary algebra but is not the displayed definition. Even dividing by \(A_n^2\) leaves the off-diagonal and tangent terms. Solving \(G\beta=b\) enforces only bubble orthogonality. A complete revision must either project onto the joint bubble-and-tangent span, or project the bubble directions within the tangent-orthogonal subspace, and then rederive the multiplier cancellation and coefficient estimates. That full repair has not been verified here. The source normalization and formulas were checked directly on the rendered pages.\par
Problem 3 (The last unweighted remainder estimate fails in dimensions at least twelve).  The last display of the proof of Lemma 6.5, at the top of p. 50, does not follow from Proposition 6.1 in the full claimed dimensional range.\par
\subsection*{Exponent and sharpness check}
The source bounds the weighted integral through the global energy norm and then asserts \[\|\nabla\rho_0\|_2^p+\|\nabla\rho_1\|_2^p
 \lesssim o(Q)+\|f\|_{H^{-1}}^p.\] Proposition 6.1 gives \(\|\nabla\rho_0\|_2\lesssim\zeta_n(Q)\). For \(n\ge7\) this only supplies \[\zeta_n(Q)^p=Q^{p^2/2}.\] It is \(o(Q)\) precisely when \(p^2/2>1\). The integer-dimensional threshold is \(7\le n\le11\); at \(n=12\), \[p=\frac75,\qquad \frac{p^2}{2}=\frac{49}{50}<1,\qquad
 \frac{Q^{49/50}}{Q}=Q^{-1/50}\longrightarrow\infty.\] The exponent decreases thereafter. For \(n=6\), the corresponding \(Q^2|\log Q|\) is indeed \(o(Q)\).\par
This is not cured by claiming an unstated improvement of the global norm. The two-bubble construction in Section 7 has \(\|f\|_{H^{-1}}\lesssim Q\) by (7.3), p. 50. The displayed lower bound on p. 52, before replacing \(Q\) by the residual, gives \(\|\nabla\rho_0\|_2\gtrsim\zeta_n(Q)\) after fixing \(L\) and increasing \(R\). In dimension twelve its \(p\)th power is bounded below by a multiple of \(Q^{49/50}\), whereas both \(o(Q)\) and \(O(Q^p)\) are \(o(Q)\). Thus the stated unweighted intermediate estimate is incompatible with that sharpness bound. The weighted integral in Lemma 6.5 can still have the required order.\par
\subsection*{A verified local replacement}
Equation reference: eq:weightedfix. Keep the weight instead of discarding it. For the affected range \(n\ge12\), Proposition 5.4, p. 44, the bound preceding (5.39), p. 45, and the integral estimate on p. 43 from Lemma B.1 give \[|\rho_0|\lesssim W,\qquad
 W^p\lesssim R^{-4(p-1)}V,\qquad
 \int VU_k\lesssim Q,\qquad |Z_k^{n+1}|\lesssim U_k.\] Consequently, \[
 \int |\rho_0|^p|Z_k^{n+1}|
 \lesssim R^{-4(p-1)}\int VU_k
 \lesssim R^{-4(p-1)}Q=o(Q).\] The remaining component is bounded, conditional on Proposition 6.4, by \[\int|\rho_1|^p|Z_k^{n+1}|
 \lesssim\|\nabla\rho_1\|_2^p
 \lesssim Q^{2p}+\|f\|_{H^{-1}}^p.\] Since \(p>1\), \(Q^{2p}=o(Q)\) and the last term is \(O(\|f\|_{H^{-1}})\) for small residual. Together with \(|\rho_0+\rho_1|^p\le2^{p-1}(|\rho_0|^p+|\rho_1|^p)\), this restores precisely the weighted estimate needed. It does not fix the decomposition in Problem 2 or certify every upstream estimate. The affected dependency is Lemma 6.5 to the interaction bound in Lemma 2.3 and then Theorem 1.4. The short available repair does not change that core dependency.\par
\section*{Overall conclusion}
The three findings distinguish a false unrestricted corollary from invalid arguments for a local theorem. The small-residual rates and the conceptual two-bubble sharpness mechanism are not refuted by the counterexample. Full verification of the repaired main theorem requires a corrected spectral decomposition and its downstream estimates, together with checking the remaining pointwise analysis. No empirical reruns are required for the analytic findings made here.\par
Finding | Issue assessment | Repair verification | Adoption / remaining work\par
Problem 1 | Major; exact counterexample | Globalization checked conditional on local theorem | Unknown; amend domain or rate\par
Problem 2 | Major; invalid core proof step | Full repair not verified | Unknown; redo projection and estimates\par
Problem 3 | Major; invalid dimensional inference | Weighted step checked conditionally | Unknown; integrate with corrected proof\par
Scientific adjudication remains pending. No human review or author adoption is recorded. A bounded current search of the author/repository and publisher records did not locate an explicit correction to these exact passages; this is not an absence or priority claim. Conclusions about the unexamined journal-final text remain open.\par
\section*{Supporting checks}
All contradictions and repairs above are analytic. An independently written arithmetic diagnostic records \(A\approx84.52127630\), \(c\approx1.01169457737\), \(c^2\approx1.02352591787\) and distance \(1/c\approx0.98844060487\), and checks rational exponents for dimensions 7–20. Floating-point rounding in the residual does not replace the exact identity [eq:residual] (\#eq:residual). No author code was executed.\par
The complete coverage ledger distinguishes read passages, conditional dependencies and unexamined sections. They are not independent acceptance or human reviewers, and no alternative model identity is asserted. Actual compilation, dependency, PDF and page-image records are retained separately from this narrative.\par
9 B. Deng, L. Sun, and J.-C. Wei, Sharp quantitative estimates of Struwe’s decomposition, author manuscript dated January 3, 2024, 59 pp. https://personal.math.ubc.ca/\textasciitilde{}jcwei/DukeFinal-2024-01-03.pdf. Associated journal identifier: https://doi.org/10.1215/00127094-2024-0026. B. Deng, L. Sun, and J. Wei, Sharp quantitative estimates of Struwe’s Decomposition, arXiv:2103.15360v2 (April 24, 2021). Earlier PDF and static TeX retained and checked only for version and relevant historical statement correspondence. https://arxiv.org/abs/2103.15360v2.\par
\end{document}
